\section{Flujo de trabajo del equipo}

El avance de las tareas se controla con un tablero ágil en \textbf{Jira Software}:

\subsection{Estados del tablero}

\begin{enumerate}[nosep]
    \item \textbf{Backlog:} Lista de requerimientos pendientes ordenados por prioridad.
    \item \textbf{To Do (Por hacer):} Tareas elegidas para el sprint que ya tienen criterios de aceptación claros (\textit{Dado / Cuando / Entonces}).
    \item \textbf{In Progress (En desarrollo):} Tarea en construcción activa dentro de una rama de Git (\texttt{feature/SIGC-XX}).
    \item \textbf{In Review / QA (En revisión):} Pull Request abierto en GitHub para que otro compañero revise el código y las pruebas.
    \item \textbf{Done (Hecho):} Tarea terminada, probada y fusionada en la rama principal (\texttt{main}).
\end{enumerate}

\subsection{Límites de Trabajo en Progreso (Límites WIP)}

Para que el equipo termine tareas antes de abrir otras nuevas, se definieron dos límites fijos:
\begin{itemize}[nosep]
    \item \textbf{En desarrollo (\textit{In Progress}):} Máximo \textbf{2 tareas al mismo tiempo} para los 3 integrantes.
    \item \textbf{En revisión (\textit{In Review}):} Máximo \textbf{2 tareas al mismo tiempo}.
\end{itemize}
\textbf{Regla del flujo:} Si la columna de revisión llega a 2 tareas, se prohíbe jalar tareas nuevas desde \textit{To Do}. Todo el equipo apoya en revisar y aprobar para liberar la columna.

\subsection{Política de bloqueo}

Si una tarea no puede continuar por un problema técnico o duda del requerimiento:
\begin{itemize}[nosep]
    \item Se marca en Jira con bandera roja (\texttt{Flagged}) y se anota el motivo en un comentario.
    \item La tarea sigue ocupando espacio en el límite WIP para que el problema no se ignore.
    \item Si el bloqueo dura más de 24 horas, el equipo se reúne para resolverlo de inmediato.
\end{itemize}

\subsection{Definición de hecho (Definition of Done - DoD)}

Una tarea solo se pasa a \textbf{Done} si cumple lo siguiente:
\begin{enumerate}[nosep]
    \item Cumple los criterios de aceptación acordados.
    \item Las migraciones de MySQL corren sin fallar (\texttt{php artisan migrate:fresh}).
    \item Las pruebas automatizadas con PHPUnit pasan en verde.
    \item El Pull Request fue revisado y aprobado por otro integrante en GitHub.
    \item El cambio está integrado en la rama principal (\texttt{main}) y probado localmente.
\end{enumerate}
